\section{Introduction}
\label{submission:sec:intro}

The proliferation of task-specific, fine-tuned deep neural networks has created a critical challenge for edge computing deployment. While parameter-efficient fine-tuning (PEFT) techniques, particularly Low-Rank Adaptation (LoRA) \cite{hu2021lora}, allow the creation of highly specialized adapters over a shared base model backbone, serving hundreds of these experts simultaneously on resource-constrained edge hardware is severely bottlenecked by memory capacity and latency constraints. Standard multi-tenant serving frameworks (e.g., S-LoRA \cite{sheng2023slora} and Punica \cite{chen2023punica}) require dynamic micro-batching and kernel-level context switching, which introduce non-negligible scheduling overheads and high hardware requirements.

To circumvent these serving overheads, recent literature has shifted toward weight-space and activation-space model merging. Rather than maintaining separate forward passes or sequential scheduling, dynamic model merging aims to blend the parameter spaces or activation streams of multi-task adapters at runtime based on the input query \cite{ilharco2022editing, wortsman2022model, yadav2023ties, stoica2023zipit}. In particular, frameworks such as Single-Pass Sample-Wise Routing with Zero-Shot Centroid Alignment (SPS-ZCA) \cite{pfsr2025} and SABLE \cite{mbh2025} perform sample-wise routing by extracting early-stage activations, projecting them onto pre-computed task centroids to extract similarity coordinates, and using these coordinates to blend downstream low-rank expert adapters in a single forward pass.

A critical, yet under-explored component of these dynamic merging architectures is the out-of-distribution (OOD) task rejection module. In a real-world multi-task serving environment, the router is inevitably subjected to queries that lie outside the set of registered tasks. Routing an OOD query to an arbitrary task adapter, or applying uniform blending weights across incompatible experts, causes catastrophic representation corruption and severe degradation of classification confidence. To avoid this routing collapse, prior works fit a diagonal Gaussian Mixture Model (GMM) in the low-dimensional coordinate space during a one-time offline calibration phase, using the log-likelihood of the similarity coordinate as a threshold for OOD rejection \cite{pfsr2025}.

In this work, we adopt the perspective of {\em The Methodologist} to critically deconstruct the implicit assumptions and potential failure modes of these coordinate density estimators. We challenge two pervasive practices in the current literature:
\begin{enumerate}
    \item {\bf The Clean Sandbox Confounder:} OOD routers are evaluated almost exclusively on clean, noise-free, and highly distinct datasets. This ignores realistic deployment environments where representation drift—stemming from sensor noise, lighting variations, or mild input perturbations—is ubiquitous.
    \item {\bf Low-Resource Calibration Overfitting:} To preserve data efficiency and minimize on-device preparation times, diagonal GMMs are fit on tiny calibration sets ($N \le 64$ samples). We hypothesize that unregularized maximum likelihood estimates of GMM parameters overfit these small samples, yielding unstable, underestimated variance boundaries that collapse under mild covariate shifts.
\end{enumerate}

To expose and rectify these limitations, we conduct a systematic sample complexity and robustness audit of coordinate density models. Crucially, we identify a major experimental design error in prior OOD evaluation pipelines: the {\em unequal noise confounder}, where representation-level noise is applied only to in-distribution samples while leaving out-of-distribution samples clean and unperturbed. This asymmetry inflates the features of ID samples and artificially collapses their similarity coordinates toward zero, completely breaking the classifier and forcing the area under the ROC curve (AUC) below $0.50$. We resolve this evaluation design flaw by introducing a mathematically sound, symmetric noise setup where identical representation-level noise is applied to both ID and OOD samples equally.

As a mathematically principled remedy for variance collapse, we introduce {\bf SRC-DE} (Shrinkage-Regularized Coordinate Density Estimation). Instead of resorting to heuristic L2 (Ridge) covariance regularizers which are non-adaptive and heavily under-regularize low-resource regimes, SRC-DE analytically estimates the optimal Ledoit-Wolf-style covariance shrinkage intensity $\alpha_{\text{opt}}$ for each mixture component. This adaptive shrinkage shrinks the unstable coordinate variances toward a stable, spherical target, preventing variance collapse and yielding robust log-likelihood boundaries without requiring hyperparameter tuning.

Our primary contributions are as follows:
\begin{itemize}
    \item We perform a comprehensive methodological audit of coordinate-space density models, exposing the hidden vulnerability of unregularized GMMs to representation-level covariate shift under low-resource calibration splits.
    \item We identify and resolve the unequal noise confounder, establishing a symmetric, scientifically valid evaluation protocol for OOD task rejection under representation-level covariate shift.
    \item We present SRC-DE, a training-free, mathematically rigorous, and adaptive covariance-shrunk coordinate density estimator designed for stable OOD rejection in dynamic model merging.
    \item We demonstrate that our proposed covariance shrinkage (SRC-DE) drastically stabilizes density boundaries under small calibration splits and covariate shift, achieving superior task rejection performance across all noise levels and calibration sizes (e.g., raising average OOD AUC under $N=16$ from $0.7701 \pm 0.0412$ to {\bf 0.7874 $\pm$ 0.0280}). In controlled scaling simulations, we show that SRC-DE consistently and globally improves OOD rejection performance under high-dimensional task registries ($K \ge 16$), yielding up to {\bf +4.63\% absolute AUC improvement} over unregularized models. We also document and resolve a critical implementation bug within standard software libraries (\texttt{scikit-learn}) concerning covariance regularization.
\end{itemize}

Furthermore, while our empirical evaluations are established using Vision Transformer backbones under image classification, we emphasize that the mathematical formulations of similarity coordinates and responsibility-weighted covariance shrinkage are completely modality-agnostic and architecture-independent. Our proposed coordinate-space density audit framework generalizes seamlessly to other major architectures, including convolutional models (via Global Average Pooling over middle feature maps) and Large Language Models (via sequence pooling of early decoder layers for text specialists prompt routing). A comprehensive, qualitative discussion outlining these multi-architecture and multi-modal generalization pathways is provided in Appendix~\ref{sec:backbone_generalization}.
